\section{Conclusion}

This paper addressed the critical challenge of fronthaul capacity limitations in uplink Cell-Free Massive MIMO systems by proposing a Task-Oriented Resource-Inference Graph Neural Network (IB-CellFree-GNN). Unlike conventional approaches that treat quantization as a signal reconstruction problem, this work formulated bit-width allocation as a dynamic, task-oriented resource optimization problem. The proposed architecture introduced a split-inference framework where distributed Access Points utilize a lightweight policy network to autonomously determine the optimal quantization precision based on local channel conditions, while the Central Processing Unit employs a Bit-Aware GNN to aggregate these multi-resolution signals. By leveraging Gumbel-Softmax relaxation, the system enabled end-to-end differentiable optimization of discrete bit-widths, allowing the network to learn a policy that prioritizes informative signal components over noise.

Numerical results confirmed that the proposed method achieves detection performance comparable to ideal full-precision systems while consuming an average of only 2.5 bits per symbol, representing a substantial reduction in fronthaul overhead. The analysis demonstrated that the system is highly robust, maintaining viable link performance even when 50\% of the Access Points are disconnected, and scales effectively in high-interference scenarios with heavy user loads. These findings validate the efficacy of integrating semantic communication principles with graph-based processing to overcome hardware constraints in distributed massive MIMO networks. Future research directions include extending this framework to joint uplink-downlink optimization to enable task-oriented precoding, and validating the inference latency and energy consumption on hardware-in-the-loop testbeds.